\chapter{Shoulder Stabilization Surgery}

\section*{Introduction \& Clinical Indications}
Shoulder stabilization surgery is a surgical intervention aimed at restoring the stability of the glenohumeral joint, which is the primary ball-and-socket joint of the shoulder. This procedure is particularly indicated for patients suffering from recurrent shoulder dislocations or persistent symptomatic subluxations, conditions where the humeral head repeatedly displaces from the glenoid cavity. The surgery typically involves repairing or reconstructing damaged soft tissue structures, such as the labrum and joint capsule, and may also address bony deficiencies of the glenoid or humeral head. The significance of this surgery lies in its ability to prevent further episodes of instability, alleviate pain, and enable patients—especially active individuals and athletes—to return to their desired functional levels and activities without fear of re-injury.

\begin{itemize}
  \item Bankart Repair
  \item Latarjet Procedure
  \item Capsulolabral Reconstruction
  \item Arthroscopic Shoulder Stabilization
  \item Open Shoulder Stabilization
  \item Remplissage Procedure
\end{itemize}

The primary surgical principle of shoulder stabilization surgery is to restore the integrity of both static and dynamic stabilizers of the glenohumeral joint, thereby preventing recurrent dislocation of the humeral head. In cases of anterior instability, the most common pathology involves a Bankart lesion, characterized by an avulsion of the anterior-inferior labrum and joint capsule from the glenoid rim. This detachment compromises the labrum's bumper effect and loosens the capsule, allowing the humeral head to dislocate.

The surgical approach aims to re-establish the anatomical relationships and tension of these structures. For soft tissue lesions, this typically involves reattaching the torn labrum and capsule securely to the glenoid rim using suture anchors. In cases of significant bony deficiencies, such as a bony Bankart lesion or a large Hill-Sachs defect, soft tissue repair alone may be inadequate. In such instances, the Latarjet procedure is employed, which involves transferring a piece of bone (the coracoid process) to augment the anterior glenoid, effectively creating a bony block and a dynamic sling effect from the transferred muscles. Other techniques, such as capsular plication or thermal shrinkage, may be utilized to tighten a redundant joint capsule, while a Remplissage procedure can address engaging Hill-Sachs lesions by tenodesing the infraspinatus tendon into the defect. The ultimate goal is to create a stable, well-centered joint capable of withstanding physiological stresses without dislocating, while preserving a functional range of motion.

The evolution of shoulder stabilization surgery has been marked by significant milestones. In 1923, British orthopedic surgeon Arthur Bankart described the essential lesion of recurrent anterior shoulder dislocation, emphasizing the importance of reattaching the avulsed labrum. His work laid the foundation for modern soft tissue stabilization techniques.

In 1954, French surgeon Michel Latarjet introduced a procedure involving the transfer of the coracoid process with its attached conjoined tendon to the anterior glenoid neck. This technique addressed cases with significant glenoid bone loss, providing both a bony block and a dynamic muscular sling effect, and remains a cornerstone for complex instability today. The advent of arthroscopy in the 1980s revolutionized shoulder surgery, allowing for minimally invasive repair of labral tears and capsular lesions using suture anchors. By the late 1990s, arthroscopic Bankart repair had largely replaced open soft tissue repairs as the gold standard for many cases of recurrent anterior instability, offering reduced morbidity and faster recovery.

Shoulder stabilization surgery is indicated in the following clinical scenarios:

\begin{itemize}
  \item Recurrent anterior glenohumeral dislocations, particularly after failed conservative management.
  \item Persistent symptomatic glenohumeral subluxation, causing pain and apprehension.
  \item First-time shoulder dislocation in young, active individuals with significant labral tears or risk factors for recurrence.
  \item Presence of significant glenoid bone loss or a large, engaging humeral head defect contributing to instability.
  \item Multidirectional instability unresponsive to comprehensive physical therapy.
  \item Failed previous non-operative treatment, including activity modification and bracing.
  \item Associated intra-articular pathologies, such as rotator cuff tears requiring surgical intervention.
  \item Patients with high functional demands requiring a stable shoulder for work or sports.
\end{itemize}

Shoulder stabilization surgery is primarily considered a definitive surgical intervention for glenohumeral instability. It is typically pursued after a trial of conservative management has failed or is deemed unlikely to succeed, particularly for patients with recurrent dislocations or significant structural damage. In these cases, surgery is essential to restore stability and function. However, it can also serve as a secondary or salvage procedure for patients who have failed previous stabilization surgeries, requiring more complex approaches to address recurrent instability or anatomical deficiencies.

\section*{Relevant Surgical Anatomy}
The shoulder joint, or glenohumeral joint, is formed by the articulation of the humeral head with the glenoid cavity of the scapula. Key anatomical structures relevant to shoulder stabilization surgery include:

- **Glenoid Labrum**: A fibrocartilaginous rim that deepens the glenoid cavity and provides stability.
- **Joint Capsule**: A fibrous envelope surrounding the joint, providing additional stability.
- **Rotator Cuff Muscles**: Comprising the supraspinatus, infraspinatus, teres minor, and subscapularis, these muscles stabilize the shoulder during movement.
- **Bony Landmarks**: McBurney's point is not applicable here, but the anterior glenoid rim and the coracoid process are critical landmarks for procedures like the Latarjet.
- **Arterial Supply**: The anterior and posterior circumflex humeral arteries supply the shoulder joint.
- **Venous Drainage**: The corresponding veins accompany the arteries and drain into the axillary vein.
- **Nerve Relations**: The axillary nerve is particularly important as it innervates the deltoid and teres minor muscles.
- **Lymphatic Drainage**: Lymphatic vessels from the shoulder region drain into the axillary lymph nodes.

Understanding these structures is crucial for the successful execution of shoulder stabilization procedures.

\section*{Preoperative Workup \& Preparation}
A comprehensive pre-operative workup is essential to optimize surgical outcomes and ensure patient safety. 

\begin{itemize}
  \item **Comprehensive Medical Evaluation**: A detailed history and physical examination assess overall health and identify comorbidities.
  \item **Imaging Studies**: MRI, often with an arthrogram, visualizes the labrum and associated soft tissue injuries, while CT scans assess glenoid bone loss.
  \item **Medical Clearance**: Patients with significant medical conditions require clearance from their primary care physician or specialists.
  \item **Medication Review**: A complete list of medications is reviewed, with blood-thinning medications typically needing to be discontinued prior to surgery.
  \item **NPO Status**: Patients are instructed to have nothing by mouth for at least 6-8 hours before surgery.
  \item **Prophylactic Antibiotics**: Intravenous antibiotics are administered shortly before the incision to reduce infection risk.
  \item **Informed Consent**: The surgeon discusses the procedure, including risks, benefits, and expected recovery, obtaining informed consent.
  \item **Pre-operative Education**: Patients receive instructions regarding post-operative care, sling use, and pain management.
\end{itemize}

Careful consideration of contraindications and precautions is vital for patient selection and safety.

\textbf{Absolute Contraindications:}
\begin{itemize}
  \item Active infection in the shoulder joint or systemic sepsis.
  \item Uncontrolled medical comorbidities that significantly increase surgical risks.
  \item Severe glenohumeral osteoarthritis, where stabilization surgery is unlikely to be beneficial.
  \item Irreparable damage to soft tissues or bone.
\end{itemize}

\textbf{Relative Contraindications and Precautions:}
\begin{itemize}
  \item Voluntary shoulder dislocation, which may hinder compliance with rehabilitation.
  \item Significant neurological deficit affecting shoulder function.
  \item Poor patient compliance or inability to adhere to post-operative rehabilitation.
  \item Extreme capsular laxity requiring more extensive procedures.
  \item Active smoking, which can impair wound healing.
  \item Unrealistic patient expectations regarding recovery.
\end{itemize}

\section*{Surgical Technique \& Procedure}
Shoulder stabilization surgery can be performed using either an arthroscopic or open technique, depending on the specific pathology and surgeon's preference. 

The procedure typically begins with the administration of general anesthesia, often supplemented by an interscalene nerve block for effective post-operative pain control. The patient is positioned in a beach chair or lateral decubitus position for optimal access to the shoulder joint.

For an **arthroscopic Bankart repair**, the following steps are performed:

\begin{itemize}
  \item Diagnostic arthroscopy is conducted to confirm the extent of the labral tear and assess for associated injuries.
  \item The anterior-inferior glenoid rim is prepared by debriding frayed tissue and decorticating the bone.
  \item Small bioabsorbable or titanium suture anchors are inserted into the prepared glenoid rim.
  \item Sutures from these anchors are passed through the torn labrum and joint capsule, reattaching them to their anatomical position.
  \item The sutures are tied, restoring the anterior bumper mechanism.
  \item The stability of the repair is assessed by manually ranging the shoulder.
\end{itemize}

For a **Latarjet procedure**, the steps include:

\begin{itemize}
  \item Harvesting the coracoid process along with its attached conjoined tendon.
  \item Preparing the anterior glenoid neck and transferring the coracoid graft, which is fixed to the anterior glenoid rim using screws.
  \item Performing a **Remplissage procedure** if indicated, which involves tenodesing the infraspinatus tendon into a humeral head Hill-Sachs defect.
  \item Closing the arthroscopic portals or open incision with sutures or sterile strips and applying a sterile dressing.
  \item A sling, often with an abduction pillow, is used to immobilize the shoulder post-operatively.
\end{itemize}

\section*{Complications \& Risks}
As with any surgical procedure, shoulder stabilization carries potential risks and complications, categorized as general surgical risks and procedure-specific risks.

\textbf{General Surgical Risks:}
\begin{itemize}
  \item **Bleeding**: Hematoma formation at the surgical site.
  \item **Infection**: Superficial wound infection or deep joint infection.
  \item **Anesthesia-related complications**: Nausea, allergic reactions, or severe cardiac events.
  \item **Pain**: Post-operative pain managed with medication.
  \item **Blood clots**: Rare occurrences of DVT or PE.
\end{itemize}

\textbf{Procedure-Specific Risks:}
\begin{itemize}
  \item **Recurrent instability or dislocation**: Failure of the repair.
  \item **Shoulder stiffness (arthrofibrosis)**: Loss of range of motion due to scar tissue.
  \item **Nerve injury**: Damage to surrounding nerves, leading to weakness or numbness.
  \item **Vascular injury**: Rare but serious injury to major blood vessels.
  \item **Hardware-related issues**: Suture anchor pull-out or irritation.
  \item **Chondral damage**: Injury to articular cartilage during the procedure.
  \item **Complex Regional Pain Syndrome (CRPS)**: A rare chronic pain condition.
  \item **Persistent pain**: Ongoing discomfort despite successful stabilization.
  \item **Failure of bone graft incorporation**: Specific to Latarjet or other grafting procedures.
\end{itemize}

\section*{Postoperative Care \& Recovery}
Post-operative recovery following shoulder stabilization surgery involves several phases. 

Immediately after surgery, patients are monitored in the Post-Anesthesia Care Unit (PACU) for pain management and stabilization. The shoulder is typically immobilized in a sling, often with an abduction pillow, to protect the surgical repair. Most patients are discharged home the same day or the following morning, provided their pain is controlled.

During the first 1-2 weeks post-operatively, the focus is on protecting the repair and managing pain and swelling. The sling is worn continuously, except for hygiene and prescribed gentle exercises. Patients receive instructions on wound care, medication schedules, and activity restrictions. A structured physical therapy program usually begins within the first week, initially focusing on passive range of motion to prevent stiffness without stressing the repair. 

Over the next 3-6 months, rehabilitation progresses through active-assisted and active range of motion exercises, followed by strengthening exercises and proprioceptive training. Return to full activity, especially contact sports or overhead activities, is typically gradual and may take 6-12 months, depending on the type of repair and individual healing.

During shoulder stabilization surgery, the surgeon documents the extent of the labral tear, capsular damage, and any associated bony lesions. The pathology report on resected tissues typically confirms the presence of a Bankart lesion or other soft tissue injuries, providing insight into the underlying pathology that necessitated surgical intervention. This documentation is crucial for understanding the surgical outcomes and guiding post-operative care.

A normal postoperative course following shoulder stabilization surgery is characterized by a gradual resolution of pain and restoration of shoulder function. Patients typically experience a significant reduction in pre-operative pain and apprehension associated with shoulder movement. The timeline for recovery varies, but most patients can expect to regain a functional range of motion within a few months, allowing them to perform daily activities without limitation. 

The repaired labrum and capsule should heal securely to the glenoid, as evidenced by clinical stability during follow-up assessments. Patients are encouraged to adhere to their rehabilitation protocol to optimize outcomes and minimize the risk of complications.

Post-operative care for shoulder stabilization surgery involves a structured follow-up plan to monitor healing and guide rehabilitation.

\begin{itemize}
  \item **Clinic Visits**: Regular post-operative visits are scheduled at 1-2 weeks, 6 weeks, 3 months, 6 months, and 1 year to assess wound healing and shoulder stability.
  \item **Suture Removal**: External sutures or staples are typically removed at the 10-14 day post-operative visit.
  \item **Physical Therapy**: A comprehensive rehabilitation program lasting 3 to 6 months is crucial for recovery, including passive, active-assisted, and active range of motion exercises.
  \item **Imaging**: Routine post-operative imaging is generally not required unless there are concerns about recurrent instability or hardware complications.
  \item **Activity Resumption**: Gradual resumption of daily activities is encouraged, with specific restrictions initially on lifting and overhead movements.
  \item **Warning Signs**: Patients are educated on warning signs to contact their surgeon, including fever, increasing redness or swelling, severe pain, or any sensation of the shoulder "giving way."
\end{itemize}

\section*{Outcomes \& Prognosis}
Shoulder instability, particularly anterior glenohumeral dislocation, is a common orthopedic injury, with an incidence ranging from 1.7 to 24 per 100,000 person-years. It predominantly affects young, active individuals, especially males aged 20-30 years, often related to sports injuries or trauma. Recurrence rates after a primary dislocation are notably high, particularly in younger patients, with rates ranging from 70-90\% in those under 20 years old. Surgical stabilization is frequently indicated for these recurrent cases, making it one of the most common shoulder procedures. Glenoid bone loss is present in up to 22\% of primary dislocations and can be found in 86\% of recurrent dislocations. Mortality associated with elective shoulder stabilization surgery is exceedingly rare, while morbidity is primarily related to the risk of recurrent instability (5-15\%), post-operative stiffness (5-10\%), and nerve injury (1-5\%).

The outcomes of shoulder stabilization surgery are generally excellent, with high rates of success in restoring stability and function. For arthroscopic Bankart repair, success rates in preventing recurrent dislocations typically range from 85-95\%, although these rates can be slightly lower in high-risk groups such as contact athletes or those with significant bone loss. Procedures addressing bone loss, such as the Latarjet procedure, demonstrate even higher stability rates, with recurrence rates reported as low as 2-5\%. 

Functional outcomes are overwhelmingly positive, with the majority of patients achieving a full or near-full return to their pre-injury activity levels, including competitive sports. Prognostic factors influencing outcomes include patient age, activity level, the extent of glenoid or humeral head bone loss, the number of prior dislocations, and adherence to the post-operative rehabilitation protocol. While recurrence remains the primary long-term concern, the overall prognosis for patients undergoing shoulder stabilization surgery is favorable, leading to significant improvements in quality of life and functional capacity.

\section*{References}
\begin{enumerate}
  \item MedlinePlus. Shoulder stabilization surgery. U.S. National Library of Medicine; 2025.
  \item Rockwood and Green's Fractures in Adults. 9th ed. Wolters Kluwer; 2020.
  \item American Academy of Orthopaedic Surgeons (AAOS). Management of Glenohumeral Instability. Clinical Practice Guideline; 2020.
  \item Netter's Atlas of Human Anatomy. 8th ed. Elsevier; 2023.
\end{enumerate}

\clearpage
